% TOP_PROBLEM: {"id": 3324, "title": "Are different notions of the crossing number the same?", "category_id": 3, "category": {"id": 3, "name": "graph_theory", "display_name": "Graph Theory", "description": "Problems involving graphs, networks, and their properties.", "slug": "graph-theory", "order_index": 3, "created_at": "2026-07-31T15:26:25.671Z"}, "status": "open", "problem_number": "OPG-37117", "set_id": 12}
% FIRST_POSTED: 2026-10-02
% ATTEMPT_STATUS: unresolved
% MODEL: GPT 6 Astra Ultra
% UnsolvedMath ID 3324; source catalogue code OPG-37117
% !TeX program = lualatex
% Research attempt dated 2026-10-02; canonical statement preserved.
\documentclass[11pt,a4paper]{article}
\usepackage[margin=25mm]{geometry}
\usepackage{fontspec}
\setmainfont{lmroman10-regular.otf}[BoldFont=lmroman10-bold.otf,
 ItalicFont=lmroman10-italic.otf,BoldItalicFont=lmroman10-bolditalic.otf]
\usepackage{amsmath,amssymb,mathtools}
\usepackage{unicode-math}
\setmathfont{latinmodern-math.otf}
\AtBeginDocument{\let\setminus\smallsetminus}
\usepackage{xurl,hyperref}
\hypersetup{colorlinks=true,urlcolor=blue}
\setcounter{secnumdepth}{0}
\setlength{\parindent}{0pt}
\setlength{\parskip}{5pt}
\setlength{\emergencystretch}{3em}
\begin{document}
\section*{Are different notions of the crossing number the same?}\label{3324}
{\small UnsolvedMath ID 3324 · OPG-37117 · Graph Theory · OpenGarden}
\subsection{Definitions and mathematical statement}
Let $G$ be a finite simple graph. Consider topological plane drawings
with distinct vertex points and simple edge arcs avoiding all other
vertices. Interior intersections are finitely many transverse
crossings; no three edges cross at one point and no two edge arcs
overlap along a segment. Two edges are allowed to cross several
times. A common endpoint is not a crossing.

For a drawing $\mathcal D$, let $c(\mathcal D)$ be its number of
crossing points, and let $p(\mathcal D)$ be the number of unordered
pairs of distinct edges whose interiors cross at least once. Pairs
are counted once even if they cross many times; this includes pairs
of adjacent edges if they cross away from their common endpoint.
Define
\[
 \operatorname{cr}(G)=\min_{\mathcal D}c(\mathcal D),
 \qquad
 \operatorname{pcr}(G)=\min_{\mathcal D}p(\mathcal D).
\]
The drawings realizing the two minima may be different.

The problem asks whether $\operatorname{pcr}(G)=\operatorname{cr}(G)$
for every finite simple graph. One always has
$\operatorname{pcr}(G)\le\operatorname{cr}(G)$ because a crossing edge
pair accounts for at least one crossing point. The difficult issue
is whether allowing repeated crossings of a pair can lower the
number of different pairs that must cross.

In defining the pair parameter, one must not restrict in advance
to drawings where each edge pair crosses at most once: that would
remove precisely the phenomenon under investigation. This parameter
also differs from odd-crossing number, which counts only pairs
crossing an odd number of times.
\subsection{Short English statement}
Can minimizing the number of crossing edge pairs ever give a smaller value than minimizing the total number of crossing points?
\subsection{Research attempt}
\paragraph{Scope and method.}
This research attempt by GPT-6 Astra Ultra is dated 2 October 2026.
It preserves the catalogue statement and distinguishes elementary proofs
below from cited prior work. No novelty claim or formal proof-assistant
verification is made. The final paragraph states the precise remaining scope.
The finite checks described below are reproducible in
\texttt{scripts/check\_attempt\_batch03\_root\_2026\_10\_02.py}.
\paragraph{Two minima and an elementary common lower bound.}
In every drawing the number of crossing pairs is at most the
number of crossing points, giving
$\operatorname{pcr}(G)\le\operatorname{cr}(G)$ after minimization.
For a drawing $D$, form an auxiliary graph whose vertices are
the edges of $G$ and whose edges are precisely the pairs that
cross in $D$. Deleting from $G$ a vertex cover of this auxiliary
graph removes every crossing in $D$. If $\mu(G)$ is the minimum
number of edges whose deletion makes $G$ planar, it follows that
\[
 \mu(G)\le\tau(\text{auxiliary graph of }D)
 \le\#\{\text{crossing pairs in }D\}.
\]
Thus $\mu(G)\le\operatorname{pcr}(G)\le\operatorname{cr}(G)$.
In particular, any graph for which $\mu(G)=\operatorname{cr}(G)$
satisfies the desired equality. This certificate uses edge
planarization, not an assumption that repeated crossings can
always be removed harmlessly.

\paragraph{Exact additivity at a cut vertex.}
Suppose $G$ is obtained from graphs $G_1,\ldots,G_s$ by
identifying one selected vertex from each, with no other
identifications or extra edges. Restrict any drawing of $G$
to each $G_i$. Crossing points internal to different pieces
are distinct, and the sets of crossing edge pairs internal
to different pieces are disjoint. Crossings between pieces
can only add to either count. Therefore each of the two
parameters of $G$ is at least the sum of that parameter over
the pieces.
For the reverse inequalities, draw each piece optimally and
make a face incident with its selected vertex the outer face.
This is possible by viewing the drawing on the sphere and
choosing the point at infinity in that face. Compress each
drawing into a separate small sector incident with the common
vertex. Their interiors are disjoint, so no crossings between
pieces are introduced. Both sums are attained. Hence
\[
 \operatorname{cr}(G)=\sum_i\operatorname{cr}(G_i),\qquad
 \operatorname{pcr}(G)=\sum_i\operatorname{pcr}(G_i).
\]
The argument can be repeated when pieces are attached one
at a time along a tree of cut vertices.

\paragraph{An infinite equality class.}
For a planar graph both parameters vanish. If a graph has
crossing number one, it is nonplanar, so its pair-crossing
number cannot be zero; the basic inequality then makes both
parameters one. It follows from additivity that every graph
assembled at cut vertices from planar pieces and pieces of
crossing number one satisfies equality, with value equal to
the number of nonplanar pieces. Copies of $K_5$ and $K_{3,3}$
give examples with arbitrarily large common value. Their
nonplanarity follows from the planar edge bounds, and explicit
drawings with one crossing prove the needed upper bounds.
Thus the class is not confined to graphs with a small total
number of crossings.

\paragraph{Why local uncrossing is not enough in general.}
Rerouting one edge beside part of another can reduce repeated
intersections of those two edges but change which third edges
it meets. If a third edge already meets the guiding edge
elsewhere, a copied segment can introduce a new crossing pair
without removing that existing pair. A proof that counts
only intersection points during this operation therefore
does not establish monotonicity of the pair count.
Karl and Toth [S3] prove the general upper estimate
$\operatorname{cr}(G)=O(\operatorname{pcr}(G)^{3/2}
\log\operatorname{pcr}(G))$ for growing pair-crossing number.
That cited comparison is much weaker than equality and is
not independently reproved here.

\paragraph{Verification and final status.}
The root script verifies exact segment drawings witnessing
one crossing for $K_5$ and $K_{3,3}$. The cut-vertex arguments
prove the asserted infinite class without extrapolating
from those checks. The missing case is general interaction
inside a piece that cannot be split at a cut vertex;
no equality theorem for all graphs is established here.

\subsection{Sources}
\begin{itemize}
\item[S1] ULAM.AI and the UnsolvedMath Contributors, supplied problems.json, record 3324 (OPG-37117), OpenGarden collection. Catalogue identity and source transcription; CC BY 4.0.
\url{https://huggingface.co/datasets/ulamai/UnsolvedMath}.
\item[S2] Open Problem Garden, Are different notions of the crossing number the same?. Original problem and discussion; the mathematical formulation below is expanded from the supplied source record.
\url{http://www.openproblemgarden.org/op/are_different_notions_of_the_crossing_number_the_same}.
\item[S3] J. Karl and G. Tóth, A slightly better bound on the crossing number in terms of the pair-crossing number (2021), comparison of the two drawing parameters.
\url{https://arxiv.org/abs/2105.14319}.

\end{itemize}
\end{document}
